\chapter{Glossar}\label{app:glossary}

%% INTERN-272: KA-134
Dieses Glossar fasst die in der Arbeit verwendeten Abkürzungen und zentralen Begriffe zusammen. Jedes
Fachbegriff-Lemma nennt zuerst den englischen Musterbegriff, dann die Definition und zuletzt die in
dieser Arbeit gleichbedeutend gebrauchten Aliase; Registry- und Code-Bezeichner stehen in
Schreibmaschinenschrift. Die Verweise zeigen auf den Abschnitt, der den Begriff einführt.

\begin{description}
  \item[ABI] Application Binary Interface --- die binär-stabile Schnittstelle, über die der Builder
  die Permutations-Module (Lebewesen) lädt.
%% INTERN-272: KA-135
  \item[Abschaltungsstufe] (\emph{instrumentation shut-off level}) Stufung, in der die
  Messeinrichtung einer Tier-Binary zurückgenommen wird: Stufe~1 schaltet in der Makro- und der
  Mikro-Ebene die Erhebung der Achsen-Parameter ab, sodass die Wall-Clock-Messung für
  \texttt{measure} und \texttt{compare} bestehen bleibt; Stufe~2 (permutativ, \texttt{release})
  entfernt auch die Wall-Clock-Messung und mit ihr die Mess-Flächen und -Interfaces aus dem
  Kompilat. Umsetzungsstand: Der Code enthält die Mess-Gates G2 und G3, deren Zustand ein
  Fingerprint-Glied ist, und den Release-Modus ohne Instrumentierung; die zweistufige Abschaltung
  als Permutations-Größe ist vorgesehen, aber nicht implementiert (Abbildung~\ref{fig:abschaltungsstufen},
  Abschnitt~\ref{sec:impl-system-axes}).
  \item[Achse] (\emph{axis}) Eine orthogonale Teilentscheidung (\enquote{Organ}) eines
  Suchalgorithmus. Das Framework ordnet seine Achsen drei Realms zu: achtzehn Organ-Haupt-Achsen
  (Kompositions-Achsen in fester Slot-Reihenfolge T0--T17, Registry-Etikett T00--T17) und eine
  optionale Organ-Meta-Meta-Achse \texttt{disk\_io}; drei System-Haupt-Achsen
  (\texttt{target\_isa}, \texttt{operating\_system}, \texttt{external\_utils}) mit ihrem
  Komplex-Wrapper; und den Mess-Realm mit der Mess-Tooling-Haupt-Achse
  (vgl.\ Abschnitt~\ref{sec:axis-triad}). Jede Achse ist Haupt-, Unter- oder Meta-Meta-Achse (siehe
  die Lemmata); in die \texttt{binary\_id} gehen ausschließlich die achtzehn Organ-Haupt-Achsen ein.
%% INTERN-272: KA-136
  \item[Achsen-Erfolgs-Parameter] (\emph{axis success parameter}) Die je Achsen-Interface
  erhobenen Kenngrößen, an denen ein Interface-Aufruf von Anfang bis Ende gemessen wird; das
  Messgerät ist der Mess-Checkpoint \texttt{checkpoint\_measure} (Mikro-Ebene), und ab der
  Makro-Ebene werden alle Parameter einschließlich der PMC-Zähler zusammengeführt
  (vgl.\ Abschnitt~\ref{sec:measurement-system}). Der Katalog der Parameter je Achse liegt nicht vor
  (Gegenstand künftiger Arbeit).
%% INTERN-272: KA-137
  \item[Adaptiv vs.\ dynamisch vs.\ messgetrieben] Begriffsabgrenzung für die Anpassung des
  Algorithmus-Verhaltens.\ \emph{Adaptiv}: Anpassung anhand laufzeit-beobachteter Datenmerkmale
  (z.\,B.\ adaptive Knotengrößen in ART)\@. \emph{Dynamisch}: im allgemeinen Sprachgebrauch jede zur
  Laufzeit getroffene Entscheidung; in dieser Arbeit enger die Laufzeit-Auswahl des fertig
  kompilierten Permutations-\emph{Moduls} über die Binary-Identität \texttt{binary\_id}
  (Modul-Ladung) --- alle achtzehn Organ-Achsen sind je Binary zur Übersetzungszeit festgelegt; hinzu kommen
  laufzeit-\emph{variierbare} Parameter dynamischer Unter-Achsen (\emph{sub-axes}: Prefetch-Distanz,
  Stapelgröße, Inline-Schwelle, Pool-Budget; vgl.\ Abschnitt~\ref{sec:measurement-system}), keine
  Laufzeit-Wahl einzelner Achsen-Strategien.
  \emph{Messgetrieben} (\emph{measurement-driven}): Auswahl der Konfiguration auf Basis empirisch
  erhobener Messreihen --- das Leitmotiv dieser Arbeit; der
  zugehörige Auswahl-Mechanismus ist in einem ersten Inkrement implementiert
  (\texttt{best\_binary\_selector}: Ranking und Auslieferung der besten Permutation je Metrik aus
  den Mess-CSVs), die vollautomatische, lastprofilübergreifende Auswahl ist Gegenstand des Ausblicks
  (vgl.\ Abschnitt~\ref{sec:outlook}).
  \item[Alder Lake] Intel-Mikroarchitektur mit heterogenen Performance- und Efficiency-Kernen; die
  zweite Messplattform dieser Arbeit (prod2, Intel Core i9-12900K) gehört ihr an. In der
  System-Registry ist sie die Rekombination \texttt{prod2\_alder\_lake} (CPU-Fabrikation
  GenuineIntel/6/151/2, RAM-Nennrate 4800~MT/s deklariert, CAS-Latenz nicht erhoben).
  \item[Anatomie] (\emph{anatomy}) Der \emph{Körper} eines Lebewesens: die feste Komposition seiner
  achtzehn Achsen-Organe und vor allem deren \emph{Verdrahtung} --- die generalisierten Interfaces,
  über die ein Organ die Dienste der übrigen nutzt (etwa eine Container-Gattung als Knoten-Substrat der
  Achse \texttt{node\_type} oder als Index-Organisation). Kein zweites Modell neben dem
  Suchalgorithmus, sondern dessen Innenstruktur (vgl.\ Abschnitt~\ref{sec:anatomy-metaphor}).
  \item[Antwort-Status] (\emph{answer status}) Die je Forschungs-Haupt- und -Teilfrage im Fazit
  ausgewiesene Stufe der Beantwortung: \emph{beantwortet} (ein Test oder eine Messung
  belegt die Antwort), \emph{strukturell beantwortet} (der Apparat, der die Antwort liefert, ist
  implementiert und getestet, die Messung liegt nicht vor) oder \emph{offen} (weder Struktur noch Messung
  liegen vor). Ohne Messwerte und konkrete Belege gilt eine Frage als nicht beantwortet, auch wenn
  ihre Struktur tragfähig ist (vgl.\ Abschnitt~\ref{sec:answers}).
  \item[AoSoA] Array of Structures of Arrays --- Speicher-Layout, das Datensätze in Blöcke fester
  Breite zerlegt und innerhalb eines Blocks feldweise (SoA) ablegt; es verbindet die
  Cache-Line-Dichte des SoA-Layouts mit der Lokalität des AoS-Layouts. In dieser Arbeit eine
  Ausprägung der Organ-Achse \texttt{memory\_layout}.
  \item[Arbeitsmodus] (\emph{work mode}) Die vier Modi der Ablauf-Kette, in dieser Reihenfolge:
  \texttt{build} (Bau der Tier-Binaries), \texttt{measure} (der deterministische Mess-Lauf),
  \texttt{compare} (Vergleich der Mess-Ebenen und Permutationen aus dem Messwertlager) und
  \texttt{release} (Auslieferung ohne Messfühler). Der Arbeitsmodus ist eine Unter-Achse des
  Mess-Realms (\texttt{work\_mode}), keine Eigenschaft der Binary-Identität; das Modus-Token
  \texttt{compare} ist von dem gleichlautenden Alias der Wallclock-Ebene in der Mess-Konfiguration
  zu unterscheiden. Alias: Modus, Ablauf-Kette.
  \item[ART] Adaptive Radix Tree~\cite{leis2013art}; Referenzentwurf mit adaptiven Knotengrößen
  (Node4/16/48/256).
%% INTERN-272: KA-138
  \item[Autonome Heuristik-Profil-Auswertung]
  (\emph{autonomous profile-driven heuristic evaluation}) Der vorgesehene Mechanismus, der aus den
  Messreihen je Lastprofil automatisch die beste Achsen-Konfiguration ableitet und als
  wiederverwendbare Heuristik-Richtlinie ablegt (Ziel-Mechanismus: ML-Klassifikation von
  Lastprofil-/Achsen-Features auf Konfigurationen). Erste Bausteine sind implementiert --- die
  automatische Best-Binary-Auswahl aus Messreihen (\texttt{best\_binary\_selector}) und die
  Extraktion/Ablage als XML-Lastprofil-Ergebnis ---; die ML-Klassifikation selbst ist
  Ausblick (vgl.\ Abschnitt~\ref{sec:outlook}).
  \item[AVX, AVX-512] x86-SIMD-Erweiterungen (256- bzw.\ 512-Bit).
%% INTERN-272: KA-139
  \item[Bestand] (\emph{stock}) Eine der vier Ablage-Klassen des Lagers, übersetzungsstatisch im
  Schlüssel-Schema festgeschrieben: Bestand~1 die gebauten Tier-Binaries, geschlüsselt über ihren
  SHA-512-Fingerprint, die Mess-Zelle daneben geklammert; Bestand~2 die Messwerte, geschlüsselt
  über den Fingerprint der vermessenen Binary und die Hardware-Identität der messenden Maschine, die
  Plattform-Zelle daneben geklammert und nie in den Digest fusioniert; Bestand~3 die
  Synthese-Batches mit den synthetisierten Mess-Kurven als Dateien; Bestand~4 die Lösungen je
  Anwender-XML (vgl.\ Abschnitt~\ref{sec:lager}).
  \item[binary\_id] Der Identitäts-Pfad einer Tier-Binary im Permutationsraum: die Belegung der
  achtzehn Organ-Haupt-Achsen in kanonischer Reihenfolge (achtzehn Segmente, das letzte
  \texttt{persistence\_target}). System- und Mess-Achsen stehen nie in der \texttt{binary\_id}
  (Registry-Attribut \texttt{binary\_id} mit dem Wert \texttt{never}); sie wirken über Bau-Version,
  Stempel-Zeilen und CSV-Spalten.
  \item[built\_stem] Feld der Plan-Textform (Fassung v1.1), das je Bau-Schritt den tatsächlich erzeugten
  Datei-Stamm ausweist --- erzeugt von derselben Funktion, die auch der Bau-Orchestrator beim
  Bau aufruft; Beleg für den \emph{einen} Hauptkanal (vgl.\ Abschnitt~\ref{sec:impl-pipeline}).
  \item[CacheEngineBuilder] (CEB) Das autonome Plattform-Ausmess-System der Kette und die zweite
  der vier Träger-Stufen; Wurzel des System-Realms. Der CEB empfängt vom Planer den freigegebenen
  System-Achsen-Raum, permutiert System- und Organ-Achsen zu Tier-Binaries, baut und stempelt sie
  je Plattform und führt die Sieben-Phasen-Pipeline über den \texttt{ExperimentDriver} aus. Sein
  Stempel enthält Mess- und System-Zeile, aber keine Organ-Zeile (die statische Sperre); sein
  SHA-512-Fingerprint ist Provenienz, kein Wiederverwendungs-Kriterium
  (vgl.\ Abschnitte~\ref{sec:mess-chain} und~\ref{sec:impl-pipeline}).
  \item[CEB] siehe \emph{CacheEngineBuilder}.
  \item[checkpoint\_measure] Das zentrale Messinstrument des Systems: ein Messpunkt, der zu einem
  Zeitpunkt die Eingangsparameter und die Zeit erhebt. Auf der Wall-Clock-Ebene ist er das
  parameter-leere Minimum; sind Makro- oder Mikro-Ebene aktiv, vollzieht er zusätzlich die
  Parameter-Klammer --- Zeit und Beobachter sind im Checkpoint vereint, die Verantwortung ist je
  Mess-Ebene getrennt. Die Zeilen werden in der Mess-Arena abgelegt, die Verschachtelung führt die
  Stapel-Arena. Alias: Mess-Checkpoint.
  \item[CLU] Cache-Line-Utilization --- Anteil genutzter Bytes je geladener Cache-Line.
  \item[CPI, IPC] Cycles per Instruction bzw.\ Instructions per Cycle.
  \item[CRTP] Curiously Recurring Template Pattern --- C++-Idiom für statischen Polymorphismus:
  Eine Klasse leitet von einer Klassenschablone ab, die sie selbst als Schablonen-Argument erhält,
  sodass die Basis zur Übersetzungszeit auf die Erbin zugreift~\cite{coplien1995crtp}. Im Framework
  die Bauform des einen Stempel-Vertrags aller Träger-Stufen (die Schablone \texttt{StempelBasis}
  mit der Erbin und der Träger-Stufe als Parametern) und der Achsen-Konzepte.
  \item[CSV-Baum] (\emph{CSV tree}) Die CSV-Form der Ergebnis-Mappe: kein einzelnes Kind-Dokument,
  sondern ein geordneter Baum aus CSV-Dateien --- ein Gesamtordner je Mappe, darunter ein Ordner je
  Blatt, darin die CSV-Dateien der einzelnen Unter-Achsen-Fahrten ---, immer aus der xlsx-Mappe erzeugt
  (xlsx nach csv, als Strategie derselben Factory) und als Darstellung ihr analog gespiegeltes
  Gegenstück; per XML wird die Mappe, der CSV-Baum oder beides am selben Ziel abgelegt.
  Umsetzungsstand: Die Naht schreibt die CSV-Projektion als eine flache Datei; die Baumform ist
  nicht implementiert (Abschnitt~\ref{sec:eval-dataflow}).
  \item[Dichte-Schwellen] (\emph{density thresholds}) Die drei festen Umschaltregeln des Prüflings
  PRT-ART zwischen seinen vier Seitendarstellungen nach Belegungsdichte, in der Reihenfolge
  Array$_{256}$ (bis 25~\%) $\to$ Array$_{65535}$ (25 bis 50~\%) $\to$ Vector$\langle$u8,u8$\rangle$
  (50 bis 75~\%) $\to$ Vector$\langle$u16,u16$\rangle$ (über 75~\%); der Prüfling schaltet nach
  diesen festen Schwellen um; die plattform-kalibrierten Regeln liegen nicht vor, ihre Festlegung
  setzt die Messläufe voraus (Teilfrage FF3.a,
  Abschnitt~\ref{sec:rqs}; Abschnitt~\ref{sec:prtart-demo}). Die Schwelle bei 50~\% ist zugleich die
  Seitentyp-Grenze zwischen Punkt- und Kompaktseite (siehe \emph{Punktseite, Kompaktseite}), keine
  weitere Umschaltregel.
  \item[Drei-Ebenen-Modell] (\emph{three-level model}) Die drei Interface-Ebenen der einen
  Architektur: \emph{Konzept} (Ebene~1, die Tierklasse als Mess-Surface --- Map, Container, Graph,
  HeuristikAdapter), \emph{Gattung} (Ebene~2, das Tier-Interface als Template --- SearchAlgorithm,
  Set, Sequence, Adapter, View, FunctionInterfaceReroute) und \emph{Genus} (Ebene~3, die Unterart als
  Implementierung des Gattungs-Interfaces --- ART, HOT, START). Darunter liegen die
  \emph{Organe}/\emph{Achsen} (Ebene~4) als eigene Klassifikation, darüber das Reich (Ebene~0);
  Tier-Binary (Ebene~5a) und Hybrid-Tier (Ebene~5b) sind die Nachphase der Kompilation. Siehe
  Abschnitt~\ref{sec:anatomy-levels}.
  \item[dTLB] Daten-Translation-Lookaside-Buffer; cacht Daten-Adressübersetzungen.
  \item[ECDF] Empirische Verteilungsfunktion (\emph{empirical cumulative distribution function})
  --- Treppenfunktion, die jedem Wert den Anteil der Beobachtungen zuordnet, die ihn nicht
  überschreiten~\cite{wilk1968probability}; im Mess-Anhang die Darstellung der Latenz-Verteilung
  über die Konfigurationen (bei einer einzigen Konfiguration eine einstufige Treppe).
%% INTERN-272: KA-140
  \item[Explizit kodierte Nicht-Erhebung] (\emph{explicitly coded non-observation}) Grundsatz der
  Mess-Kette: Ein nicht erhobener Wert ist ein eigener Zustand (NaN nach IEEE~754 bzw.\ ein
  Gültigkeits-Flag~\cite{ieee754_2019}) und wird nie als Null oder Phantomwert geführt --- weder im
  Schema noch in CSV, xlsx, Tabelle oder Diagramm; eine Messreihe ohne gültige Zeile erzeugt keine
  Ausgabe statt einer leeren. Der fehlende Wert bleibt damit als eigene Kategorie
  auswertbar~\cite{graham2009missing}. Alias: Grundsatz der ausgewiesenen Nichterhebung; im Code das
  Vokabular \texttt{honest-0}/\texttt{honest-empty} der Kommentare und Prüf-Dock-Meldungen, im
  Diagramm-Generator des Werkzeug-Repositories (\texttt{diagram\_generator}) die Platzhalter-Routine
  \texttt{write\_honest\_empty\_placeholder} (vgl.\ Abschnitt~\ref{sec:axis-triad}).
  \item[Eytzinger-Layout] (\emph{Eytzinger layout}) Anordnung eines sortierten Feldes in der
  Ebenen-Reihenfolge eines impliziten vollständigen Binärbaums (Kinder von Index~$i$ bei $2i$ und
  $2i+1$), sodass die binäre Suche cache-line- und prefetch-freundlich
  verläuft~\cite{khuong2017eytzinger}. In dieser Arbeit eine Ausprägung der Speicher-Layout-Achse.
  \item[Fingerprint] (\emph{fingerprint}) Der kryptografische Identitäts-Hash je Träger-Stufe: Der
  Planer trägt einen SHA-256 über Version, ISA und Betriebssystem (64 Hex-Zeichen); die Tier-Binary
  einen SHA-512 über elf Glieder in fester Ordnung (Format~6: Format-Kennung, Mess-, System- und
  Organ-Zeile, Unter-Achsen-Werteset, Toolchain, Enabled-Mengen-Signatur der Bau-Achsen,
  Overlay-Quellmenge, Mess-Gates, Hybrid-Komposit, Bau-Versions-Basis; 128 Hex-Zeichen); der CEB
  einen SHA-512 als Provenienz. Der Fingerprint ist der Lager-Schlüssel des Bestands~1 und das
  Wiedererkennungs-Kriterium einer Permutation.
  \item[First-Touch] (\emph{first-touch policy}) NUMA-Platzierungsregel des Betriebssystems: Eine
  Speicherseite wird dem Knoten zugeordnet, dessen Prozessor sie zuerst beschreibt; deshalb
  entscheidet der initialisierende Thread über die Speicher-Lokalität~\cite{lameter2013numa}.
  \item[Forest-Plot] (\emph{forest plot}) Darstellung mehrerer Effektschätzer mit ihren
  Unsicherheitsintervallen als übereinander gestellte Punkt-Balken-Zeilen~\cite{lewis2001forest};
  im Mess-Anhang die Sicht auf den Achsen-Austausch je Konfiguration gegen eine Referenz.
  \item[Gattung] (\emph{family}; Code: \texttt{AnatomyGenus}) Die zweite Ebene des Frameworks
  (Ebene~2) und die erste Tier-Interface-Ebene unter einem Konzept: das Interface, das eine Gattung als
  UND-Schnittmenge der Interfaces aller ihrer Genera nach außen als Template trägt (im Metapher-Kanon
  der Elefant $=$ SearchAlgorithm unter dem Säugetier Map). Sechs Gattungen: SearchAlgorithm unter
  Map, Set, Sequence, Adapter und View unter Container (Specht, Storch, Kakadu, Falke) sowie
  FunctionInterfaceReroute (Ameise) unter HeuristikAdapter. Die Gattung ist die einzige gepflegte
  Quelle der Einordnung; das Konzept wird zur Laufzeit aus ihr ermittelt. Je andockender Gattung gibt
  es genau ein Prüf-Dock (fünf Docks für SearchAlgorithm, Set, Sequence, Adapter und View), während
  die klassifizierende Gattung FunctionInterfaceReroute keines hat und ihre Binaries am Dock der
  gemeldeten Ziel-Gattung andocken. Die Tier-Binary zeigt dem Prüf-Dock ihr Gattungs-Interface; Docks
  verschiedener Gattungen werden nie gemischt. Die Gattungs-Tests des Konformitäts-Gates sind generisch für alle
  Genera der Gattung. Code-Brücke: \texttt{AnatomyGenus} (Umbenennung im Code nach dem Vortrag; siehe
  \emph{Konzept}, \emph{Genus}, \emph{Drei-Ebenen-Modell}).
%% INTERN-272: KA-141
  \item[Genus] (\emph{genus}) Die dritte Ebene (Ebene~3) und zweite Tier-Interface-Ebene: die
  Unterart, die Implementierung des Gattungs-Interfaces mit ggf.\ erweiterten Spezial-Interfaces (im
  Metapher-Kanon die Elefanten-Arten Mammut $=$ ART, afrikanischer Elefant $=$ HOT, indischer Elefant
  $=$ START unter der Gattung SearchAlgorithm; die Giraffe ist reserviert, CoCo hat kein Code-Objekt).
  Unter dem Genus liegen eine oder mehrere Suchstrategien (einfach $=$ Belegung einer einzelnen Achse,
  komplex $=$ Komposition mehrerer Achsen); die Organe sind eine eigene Klassifikation (Ebene~4), die
  die Genus-Implementierung bindet. Das Genus ist weder die Suchstrategie-Achse noch die Komposition
  seiner Organe (die Zusammensetzung aus Achsen ist Sache des Versions-Stempels,
  Abschnitt~\ref{sec:stamp-model}, nicht der Taxonomie). Die Spezial-Interfaces des Genus prüft das
  Konformitäts-Gate nur unter \texttt{--debug} (Genus-Tests); dem Prüf-Dock zeigt die Tier-Binary ihr
  Gattungs-Interface, nicht das Genus. Die Genus-Implementierung führt je Funktion einen Hauptalgorithmus,
  der ausschließlich über Achsen-Interface-Aufrufe arbeitet, und ruft fremde Achsen als erste auf. Ist-Stand:
  Der Code trägt das Genus heute nur als Achsen-Belegung; die Genus-Implementierungsklasse fehlt
  (erkannter Architekturfehler, Behebung nach dem Vortrag). Das Hybrid-Artefakt klassifiziert die
  Gattung FunctionInterfaceReroute; \texttt{genus()} einer Hybrid-Binary meldet dagegen die geerbte
  Ziel-Gattung, die Umleitung bleibt nach außen transparent. Code-Brücke: \texttt{AnatomyGenus}
  bezeichnet im Code die \emph{Gattung} dieser Arbeit, nicht das Genus (Umbenennung nach dem Vortrag;
  siehe \emph{Gattung}).
  \item[Gitlink] (\emph{gitlink}) Der im übergeordneten Repositorium verzeichnete Commit eines
  Submoduls (Eintragsmodus 160000); \texttt{.gitmodules} deklariert je Submodul Pfad, URL und
  optional den Zweig~\cite{git2026gitmodules}. Das Super-Repositorium führt so den maßgeblichen
  Stand von Engine, Prüfling und Manuskript.
  \item[GoF] \enquote{Gang of Four} --- die kanonischen Entwurfsmuster nach Gamma et
  al.~\cite{gamma1994patterns}.
%% INTERN-272: KA-142
  \item[golden-Referenzkatalog] (\emph{golden reference catalogue}) Der feste Regressions-Detektor
  des Ganzsystems: siebzehn Organ-Achsen mit je zwei Ausprägungen, \texttt{persistence\_target} auf
  den In-Memory-Vertreter festgelegt, ergeben $2^{17} = 131072$ Identitäten, über eine CRC-64-Prüfsumme
  statt über die materialisierte Liste verankert --- unter der vereinfachenden, überholten Annahme genau
  zweier Auslegungen je Achse; die Achsen tragen mehrere Auslegungen, der Katalog bleibt Regressions-Detektor,
  nicht Obergrenze. Er ist von der Obergrenze der Binary-Identitäten je System-Permutation zu trennen: dem
  Produkt der je Achse im XML-Experiment gewählten Auslegungen, das der Planer je Profil und Plattform dynamisch
  aus den aktivierten Katalogen rechnet und das erst zur Ausführung bekannt ist
  (vgl.\ Abschnitt~\ref{sec:mess-chain}).
  \item[H2-Akte] (\emph{H2 dossier}) Die werkzeug-berechnete Datenquelle der Evaluations-Hypothese
  H2 (Code-Qualität je Quelle): eine gewichtete cppcheck-Befunddichte je kLOC über die vendorten
  Originalquellen der Stand-der-Technik-Entwürfe; ohne vendorte Originalquelle führt die Akte
  \texttt{n/a} mit Grund. Nicht zu verwechseln mit den PRT-ART-Telemetrie-Metriken H1--H3.
  \item[Haupt-Achse] (\emph{main axis}) Eine übersetzungsstatische Achse, deren Wahl fest in den
  CEB bzw.\ die Tier-Binary einkompiliert wird und in der Bau-Legende erscheint: die achtzehn
  Organ-Haupt-Achsen (identitätsbildend), die drei System-Haupt-Achsen und die
  Mess-Tooling-Haupt-Achse. Ihre Unter-Achsen permutiert der Planer zur Laufzeit. Im
  Stempel-Eintrag ist die Haupt-Achse die klammerlose Ebene~0.
  \item[HBM] High-Bandwidth Memory.
  \item[HDR-Histogramm] (\emph{high dynamic range histogram}) Histogramm mit logarithmisch
  gestuften Klassen und garantierter relativer Genauigkeit über viele
  Größenordnungen~\cite{tene_hdrhistogram}; in der Auswertung die Verteilungssicht auf die
  Latenzen, während die Perzentile nearest-rank aus den Rohwerten bestimmt werden (siehe
  \emph{nearest-rank-Median}).
  \item[header-only] Bibliotheksform, deren gesamter Code in Header-Dateien liegt und beim
  Einbinden übersetzt wird; es entsteht kein eigenes Bibliotheks-Artefakt. Der Prüfling PRT-ART
  und die Hybrid-Stufe liegen in dieser Form vor.
  \item[HOT] Height Optimized Trie~\cite{binna2018hot}.
%% INTERN-272: KA-143
  \item[Hybrid-Tier-Binary] (\emph{hybrid tier binary}) Die vierte Träger-Stufe (Ebene~5b: zwischen dem
  CEB-Prüfdock und den Tier-Binaries eingeordnet als Multiplexer mit Slots) in der Bau-Reihenfolge, in
  der Verarbeitungskette der optionale Mux-Einschub zwischen CEB-Prüf-Dock und Tier-Binary: ein Modul der
  Konzept HeuristikAdapter (Gattung FunctionInterfaceReroute), das zur Übersetzungszeit die
  Interfaces einer Gattung und einer Gattung erbt und Aufrufe nach heuristischen Kriterien an
  angeschlossene Tier-Binaries durchstellt; es exportiert dieselben Pflicht-Symbole wie ein Tier,
  meldet nach außen die Ziel-Gattung und wird ausschließlich mithilfe des CEB gebaut.
  Umsetzungsstand: Klassifikation, Dock-Schicht, XML-Parser und Binary-Proxy sind implementiert; der
  Modul-Export als eigenes Kompilat, der Reroute-Router, die Verdrängung, die
  Zustands-Synchronisation (Memento-Sidecar) und das Live-Ranking sind nicht implementiert
  (vgl.\ Abschnitt~\ref{sec:impl-hybrid-lager}).
  \item[HybridAware] Der sechste Kern-Modus des Mess-Modells und Wert \texttt{HybridAware} der
  Scheduling-Dimension \texttt{hetero\_core\_dispatch} (neben \texttt{None}, \texttt{PCoresOnly} und
  \texttt{ECoresOnly}): Der Lauf kennt die heterogene Kern-Topologie, pinnt seine Threads zur Laufzeit
  über die Affinitäts-Schnittstelle des Betriebssystems auf ausgewählte Kerne beider Sorten und misst
  dort mit dem zur Kern-Sorte gehörenden Zähler-Satz; Voraussetzung der Cache-Line-Awareness-Messung
  und deshalb Pflicht-Bestandteil des Mess-Modells (vgl.\ Abschnitt~\ref{sec:axis-triad}).
  \item[Instrumentierungs-Artefakt] (\emph{probe effect}) Die Veränderung des beobachteten
  Verhaltens durch das Mess-Instrument selbst~\cite{gait1986probe}; in dieser Arbeit der
  Unterschied zwischen der Tier-Binary mit und ohne Messfühler, der über die
  Plausibilitäts-Invariante (mit Messfühler langsamer als ohne) geprüft und über die
  Messfehler-Interpolation herausgerechnet wird.
  \item[Interface-Fläche] (\emph{interface surface}) Eine der drei Vertragsflächen, die jede
  Träger-Stufe bietet: Fläche~1 das Gattungs-Interface (Laufzeit, Abstract Factory), Fläche~2 der
  Stempel (Übersetzungszeit, ABI-stabil), Fläche~3 der Mess-Durchstich (\texttt{IMessVisitor}), der
  die Messung als eigene Naht quert, damit Konzept- und Gattungs-Interfaces unverändert bleiben
  (Abbildung~\ref{fig:traeger-flaechen}). Alias: Vertragsfläche, Träger-Fläche.
  \item[IQR] Interquartilsabstand (\emph{interquartile range}) --- Differenz zwischen drittem und
  erstem Quartil (p75 $-$ p25); im Mess-Anhang das Streuungsmaß der Geschwister-Paar-Differenzen in
  den Austauschbarkeits-Langtabellen und den Forest-Plots (Whisker = IQR); ein Box-Plot der
  Latenzen wird dort bewusst nicht gezeichnet, weil das Schema je Operation nur die Perzentile p50,
  p99 und p999, aber keine Quartilsgrenzen (p25/p75) trägt~\cite{mcgill1978boxplots}.
  \item[Komplex-Achse] (\emph{complex axis}) Klammer über einer festen Rekombination mehrerer
  Haupt-Achsen, die gemeinsam wie \emph{eine} Achse auftreten: die äußere Klammer über Ziel-ISA,
  Betriebssystem und Erweiterungs-Hardware sowie die innere Komplex-Achse \texttt{target\_isa}
  (RAM-Frequenz, CAS-Latenz, CPU-Fabrikation); im Identitäts-Stempel belegt eine Komplex-Achse
  genau ein Feld (vgl.\ Abschnitt~\ref{sec:complex-axis}). Ein nicht deklariertes Glied wird als
  nicht deklariert gestempelt, nie geraten; die CAS-Latenz ist für keine Messmaschine erhoben (n/a:
  die SMBIOS-Struktur Typ~17 führt sie nicht).
  \item[Konformitäts-Gate] (\emph{conformance gate}) Die unbedingte Prüfung vor jeder Messung: Das
  Gate prüft am Prüf-Dock den \texttt{std::map}-Vertrag der Konformitäts-Teilmenge
  (Anhang~\ref{app:interfaces}), damit nur Vergleichbares verglichen wird. Unter \texttt{--debug}
  fährt es zusätzlich den Standard-Testsatz des Tier-Interfaces: den generischen Gattungs-Teil
  (gültig für alle Genera) und den Genus-speziellen Teil (die erweiterten Spezial-Interfaces); dieser
  indirekte Debug-Modus fällt mit dem Release-Bau zusammen und wird nicht separat geführt (siehe
  \emph{Gattung}, \emph{Genus}).
  \item[Konsument] (\emph{consumer}) Die Rolle des Diplomarbeits-Codes gegenüber der Engine: Ein
  Konsument linkt ausschließlich gegen die Cache-Engine und benutzt eine Tür
  (\texttt{get\_cache\_engine()}); der Mess-Treiber ist zugleich ihr Treiber. Davon zu trennen ist
  der Bediener --- der Mensch, der Planer und Anlage steuert. Alias: formaler Konsument und Treiber.
  \item[Konzept] (\emph{concept}; Code: \texttt{AnatomyGattung}) Die oberste Ebene unter dem Reich
  (Ebene~1): die Tierklasse als Mess-Surface --- Map (Säugetier), Container (Vogel), Graph (Pflanze),
  HeuristikAdapter (Insekt) ---, das nach außen sichtbare Mess-Interface, das die Messbarkeit
  definiert; je Gattung materialisiert ein Prüf-Dock sie. Das Konzept wird nicht als zweites Datum
  gepflegt, sondern zur Laufzeit aus der Gattung ermittelt, über ein konzeptunabhängiges C-Symbol
  (das Symbol trägt für jedes Konzept denselben Namen). In der Produktlinien-Lesart ist das Konzept
  die Feature-Kern-Komposition, die Gattung die Feature-Spezial-Komposition. Lesehinweis:
  \enquote{Adapter} ohne Präfix meint stets die Container-Gattung, das Hybrid-Konzept heißt immer
  HeuristikAdapter. Code-Brücke: \texttt{AnatomyGattung}, ABI-Symbol
  \texttt{comdare\_anatomy\_gattung} (Umbenennung im Code nach dem Vortrag; siehe
  \emph{Drei-Ebenen-Modell}, \emph{Gattung}, \emph{Genus}).
%% INTERN-272: KA-144
  \item[Lager] (\emph{repository}) Der persistente Bestand des Apparats, in dem Gebautes und
  Gemessenes unter Stempel-Schlüsseln liegt (siehe \emph{Bestand}): Der Lager-Baum ist nach
  Konzept, Gattung und Realm gegliedert, eine Permutation wird über ihren Fingerprint wiedererkannt
  statt neu gebaut, die Messung ist immer eingeschaltet und misst im Lager nur nach, was dort
  fehlt. Jede Maschinen-Konfiguration hat ihren eigenen Messwert-Bestand; ein
  Algorithmus-Versionswechsel löst den Neubau des betroffenen Lager-Stapels samt Neumessung aus.
  Umsetzungsstand: Bestand~1 und~2 sowie das Schlüssel-Schema aller vier Bestände sind implementiert;
  die Mechanik von Bestand~3 und~4, die Kurven-Facette, der Ablage-Pfad des Hybrid-Tiers und die
  Zielwahl über die XML bilden den Lager-Vollausbau und sind nicht implementiert (Abbildung~\ref{fig:lager-bestaende},
  Abschnitt~\ref{sec:lager}).
  \item[Lastprofil-Selektionsverzerrung] (\emph{workload selection bias}) Die Verzerrung eines
  Vergleichs dadurch, dass die Lastprofile so gewählt sind, dass ein Entwurf unter ihnen bevorzugt
  gut abschneidet~\cite{blackburn2016truth}; ihr begegnet die Arbeit mit dem globalen
  Standard-Workload als Permutation über alle verfügbaren Workloads und mit je Veröffentlichung
  reproduzierten Profilen.
  \item[Lebewesen] (\emph{living being}) Ein vollständig komponiertes Tier-Modul eines Konzepts ---
  alle Gattungen und Genera sind Lebewesen; die Gattung SearchAlgorithm ist der Fall mit der vollen
  Achtzehn-Achsen-Anatomie, und nur die ABI-Adapter-Sicht verengt den Begriff auf ihn (technischer
  Identifier dort: SearchAlgorithm-Instanz). Die Anwender-XML führt ein Lebewesen mit einem oder
  mehreren Tieren. Abgrenzung: Achsenlose Graph-Module sind Viren (siehe \emph{Virus}).
  \item[LOUDS] Level-Ordered Unary Degree Sequence --- succincte (platzsparend nahe der
  informationstheoretischen Schranke) Baum-Kodierung, die die Knotengrade ebenenweise unär
  ablegt~\cite{jacobson1989louds}.
%% INTERN-272: KA-145
  \item[M0--M3] (\emph{meta-levels M0--M3}) Die vier Abstraktionsstufen, nach denen diese Arbeit
  ihre Achsen-Begriffe ordnet, in Anlehnung an die vierstufige Meta-Ebenen-Architektur der Object
  Management Group (MOF)~\cite{omg2016mof}: M3 die Meta-Meta-Achsen als Indirektion einer
  Beschreibung; M2 die Meta-Achsen als literale Klassifizierung der unter ihrem Begriff geführten
  Haupt- und Unter-Achsen; M1 Konzeption und Laufzeit-Implementierung; M0 der Binärcode. Die
  Zuordnung von M3 ist gesetzt, der Schnitt zwischen M2 und M1 ist Lesart
  (Abbildung~\ref{fig:m0-m3}, Abschnitt~\ref{sec:impl-system-axes}).
  \item[MESI, MOESI, MESIF] Cache-Kohärenzprotokolle (Grund-Protokoll bzw.\ AMD-/Intel-Varianten).
  \item[Mess-Ebene] (\emph{measurement level}) Eine der drei Ebenen der Mess-Tooling-Haupt-Achse,
  in genau dieser Reihenfolge: \emph{wallclock} (w) --- die Wall-Clock der Gesamtläufe im
  CacheEngineBuilder über die Last-Rahmenwerke; \emph{macro} (ma) --- die Zeit der
  Gattungs-Interface-Funktionen der Tier-Binary; \emph{micro} (mi) --- Zeit und Erfolgs-Parameter je
  Achsen-Interface. Teilmengen sind erlaubt, jede andere Reihenfolge ist syntaktisch falsch
  (vgl.\ Abschnitt~\ref{sec:measurement-system}). Alias: w/ma/mi; Wallclock-, Makro-, Mikro-Ebene.
  \item[Messfühler] (\emph{probe}) Die in eine Tier-Binary einkompilierte Messeinrichtung
  (Mess-Durchstich, Beobachter, PMC-Zähler); ihre Belegung ist im Mess-Gates-Glied des Fingerprints
  kodiert, sodass die Binary mit und die Binary ohne Messfühler verschiedene Identitäten besitzen und
  beide im Lager liegen. Der Unterschied ihrer Laufzeiten ist das Instrumentierungs-Artefakt.
  \item[Mess-Realm] (\emph{measurement realm}) Der dritte Achsen-Realm neben Organ- und
  System-Achsen: die Belange des Messens selbst, verwurzelt am Experiment-\emph{Planer}. Seine
  Haupt-Achse ist das Mess-Tooling (wallclock/macro/micro), die Auffächerung des CEB-Typs; seine
  Unter-Achsen sind die sechzehn Mess-Kategorien (CSV-Spalten), der Arbeitsmodus, Workloads und
  Datensätze sowie die Rückschrieb-Methoden; das Last-Rahmenwerk ist seine Meta-Meta-Haupt-Achse.
  Alle Mess-Realm-Achsen sind strikt binary-id-neutral (vgl.\ Abschnitt~\ref{sec:axis-triad}).
%% INTERN-272: KA-146
  \item[Meta-Meta-Achse] (\emph{meta-meta axis}) Eine Achse der Stufe M3: Sie beschreibt nicht
  eine Belegung, sondern die Indirektion einer Beschreibung, und stempelt als Klammer-Anhang ihrer
  Realm-Zeile; die Meta-Meta-Ebene eines Stempel-Eintrags ist seine Klammer-Tiefe. Je Realm gibt es
  einen eigenen Diskriminator: \texttt{organ\_meta\_meta} (\texttt{disk\_io}),
  \texttt{system\_meta\_meta} (die Familien-Achsen unter dem Hub der Erweiterungs-Hardware, etwa
  die SIMD-Familie) und \texttt{measurement\_meta\_meta} (\texttt{load\_framework}); die
  PMC-Erfassung zählt zu den Meta-Meta-Mess-Achsen.
%% INTERN-272: KA-147
  \item[Metapher-Kanon] Die biologische Lesart der Architektur, in dieser Arbeit einheitlich
  gebraucht: Reich $=$ Animalia (\texttt{IAnatomyBase}); Tierklasse $=$ Konzept (\enquote{Säugetier}
  $=$ Map, \enquote{Vogel} $=$ Container, \enquote{Pflanze} $=$ Graph --- botanisch das Reich Plantae,
  die Metapher greift hier bewusst über das Reich Animalia hinaus ---, \enquote{Insekt} $=$
  HeuristikAdapter); Art $=$ Gattung (\enquote{Elefant} $=$ SearchAlgorithm, \enquote{Specht} $=$ Set,
  \enquote{Storch} $=$ Sequence, \enquote{Kakadu} $=$ Adapter, \enquote{Falke} $=$ View,
  \enquote{Ameise} $=$ FunctionInterfaceReroute); Unterart $=$ Genus (\enquote{Mammut} $=$ ART,
  \enquote{afrikanischer Elefant} $=$ HOT, \enquote{indischer Elefant} $=$ START); Organ $=$ Achse;
  Anatomie $=$ der Körper eines Lebewesens; Tier $=$ ein Lebewesen einer Gattung, das Individuum die
  Tier-Binary; Virus $=$ achsenloses Graph-Modul; Blut $=$ das flüssige Organ, das alle
  Organ-Achsen verbindet und im Planer zu einer Übersetzungszeit-Konfiguration verdichtet wird;
  Lebensraum $=$ die System-Achsen; Labor $=$ die Mess-Achsen
  (vgl.\ Abschnitt~\ref{sec:anatomy-metaphor}). Wo der Text den Mess-Querschnitt des Hosts meint,
  sagt er das ohne Metapher.
  \item[MLP] Memory-Level Parallelism --- Anzahl gleichzeitig ausstehender Speicherzugriffe, die
  ein Kern überlappen kann; begrenzt durch die Miss-Status-Holding-Register (siehe \emph{MSHR}).
  \item[MPKI] Misses per Kilo-Instructions --- Cache- oder TLB-Fehlzugriffe je tausend
  ausgeführter Instruktionen; normiert Fehlzugriffszahlen über Programme verschiedener Länge.
  \item[MSHR] Miss Status Holding Register --- Puffer eines Caches, der ausstehende Fehlzugriffe
  verwaltet und damit die maximale Zahl gleichzeitig offener Zugriffe (MLP) begrenzt.
  \item[nearest-rank-Median] (\emph{nearest-rank median}) Perzentil-Definition ohne
  Interpolation: Das $p$-Quantil ist der Wert an Rang $\lceil p \cdot n \rceil$ der sortierten
  Stichprobe (Typ~1 nach Hyndman und Fan~\cite{hyndman1996quantiles}); alle Perzentile der
  Auswertung werden so aus den Rohwerten bestimmt, die HDR-Verteilung ergänzt die Sicht.
  \item[NINE] Non-Inclusive-Non-Exclusive --- Inklusions-Politik eines gemeinsamen
  Last-Level-Cache, der gegenüber den privaten Caches weder Inklusion noch Exklusion garantiert.
  \item[NUMA] Non-Uniform Memory Access.
  \item[OLC] Optimistic Lock Coupling --- optimistische Nebenläufigkeits-Technik (ARTSync,
  P08)~\cite{leis2016olc}.
  \item[Organ] (\emph{organ}) Metaphorischer Begriff für eine \emph{Achse} (Ebene~4 unterhalb des
  Drei-Ebenen-Modells); ein Suchalgorithmus besitzt achtzehn Organ-Achsen. Neben den achtzehn
  Organ-Haupt-Achsen führt der Organ-Realm eine optionale Organ-Meta-Meta-Achse
  (\texttt{disk\_io}), die nur bei einem Platten-Rückschreibe-Ziel wirksam wird (siehe
  \emph{Achse}, \emph{Meta-Meta-Achse}).
  \item[Overlay-Quellmenge] (\emph{overlay source set}) Die eine Menge der Quelldateien, über die
  das Overlay-Glied des Tier-Fingerprints hasht: je Achse die Dateien ihrer eigenen Implementierung
  plus die gemeinsame Tier-Substanz, in fester Ordnung je Achsen-Kategorie konkateniert und einmal
  per SHA-512 gehasht; dieselbe Menge ist die Grundgesamtheit des Versions-Locks.
  \item[Pareto-Front] (\emph{Pareto front}) Menge der Konfigurationen, die in keiner Metrik
  verbessert werden können, ohne sich in einer anderen zu verschlechtern
  (Pareto-Dominanz)~\cite{emmerich2018multiobjective}; im Mess-Anhang die Sicht auf den
  Latenz-Durchsatz-Speicher-Kompromiss der Permutationen.
%% INTERN-272: KA-148
  \item[Passungs-Stempel] \emph{Entwurf} eines Abgleichs, der die einkompilierte
  Plattform-\emph{Klasse} einer Tier-Binary bei jedem Mess-Lauf gegen das Laufzeit-Verdikt der
  Hardware-Erkennung prüft; nicht implementiert (vgl.\ Abschnitt~\ref{sec:answers}).
  \item[P-Core, E-Core] (\emph{performance core, efficiency core}) Die zwei Kern-Sorten hybrider
  Prozessoren (auf Intel seit Alder Lake); die Messung pinnt Threads je Kern-Sorte und erhebt die
  PMC-Zähler getrennt. Die Kern-Zuteilung ist die Dimension \texttt{hetero\_core\_dispatch} der
  Scheduling-Unter-Achse an \texttt{target\_isa}, die Kern-Klasse die Laufzeit-Unter-Achse
  \texttt{core\_class}; das Mess-Modell führt sechs Kern-Modi (vier Pinning-Modi je Kern-Sorte, der
  ungepinnte Referenzlauf und HybridAware als sechster Modus), die Anbindung an das Lauf-Profil ist
  nicht implementiert (vgl.\ Abschnitt~\ref{sec:axis-triad}).
  \item[Pitchfork-Layout] (\emph{Pitchfork layout}) Konvention für die Verzeichnisstruktur nativer
  C-/C++-Projekte (\texttt{apps/}, \texttt{libs/}, \texttt{tools/}, \texttt{tests/},
  \texttt{docs/}, \texttt{external/} --- in der Engine als \texttt{ext/}), der die Engine
  folgt~\cite{pike2018pitchfork}.
  \item[Planer] (\emph{planner}) Der Experiment-Planer (\texttt{ExperimentPlanDirector}): löst die
  Anwender-XML gegen die drei Angebots-Kataloge auf, wählt die Mess-Achsen und steuert die
  Builder-Aufträge der Kette; Wurzel des Mess-Realms und erste der vier Träger-Stufen, als Werkzeug
  ein Subkommando-Programm (\texttt{validate}, \texttt{plan}, \texttt{cache-key}, \texttt{fingerprint},
  \texttt{status}, \texttt{check-size}, \texttt{simulate}, \texttt{version}, \texttt{help}); sein Stempel
  trägt als einziger die Versionszeile und einen SHA-256-Fingerprint (vgl.\ Abschnitt~\ref{sec:mess-chain}).
  \item[PMC] (Hardware) Performance Monitoring Counter. Im Mess-Realm ist der Zählerpfad ein
  Baustein der Kollektor-Achse (Regime PmcCounter); die modellgebundenen Roh-Ereignisse liefert ein
  Katalog je Mikroarchitektur, und in der Achsen-Typologie zählt PMC zu den Meta-Meta-Mess-Achsen.
  \item[Provenienz-Stufe] Vertrauensstufe eines erhobenen Hardware-Werts --- konfiguriert-gemessen,
  SPD-JEDEC-Basiswert oder deklariert-nicht-gemessen ---, die jedes Ergebnis samt
  Erhebungszeitpunkt mitführt; eine niedrigere Stufe verdrängt eine höhere nur mit mitwandernder
  Kennzeichnung (vgl.\ Abschnitt~\ref{sec:hw-probe}).
%% INTERN-272: KA-149
  \item[Prüf-Dock] (\emph{probe dock}) Die builder-seitige Mess-Orchestrierung je Gattung: Das Konzept
  definiert die Messbarkeit, je Gattung materialisiert ein Prüf-Dock sie. Genau ein
  Dock je Gattung (SearchAlgorithm, Set, Sequence, Adapter, View) lädt ein Lebewesen, treibt es durch
  die Interface-Operationen und vermisst es; Hybrid-Docks binden angeschlossene Tiere. Das Dock ist
  keine ABI-Grenze --- die ABI-Grenze bleibt das gattungs-eigene Antriebs-Sub-Interface samt
  POD-Snapshot; die Persistenz der Ergebnisse übernehmen der Ergebnis-Einzug in den
  Experiment-Baum und der Rückschrieb (vgl.\ Abschnitt~\ref{sec:anatomy-levels}).
  \item[Prüfling] Der testweise einzugliedernde neue Suchalgorithmus (hier PRT-ART); wird vom
  Builder gegen die Kernbestandteile der Cache-Engine kompiliert.
  \item[PRT-ART] Probst-Redirect-Tree-/ART-Hybrid --- der Prüfling dieser Arbeit.
  \item[Punktseite, Kompaktseite] (\emph{array-indexed page, vector page}) Die zwei Seitenformen
  des PRT-ART-B$^+$-Knotens nach Belegungsdichte: Die Punktseite adressiert Einträge direkt über
  ein Feld (8- bzw.\ 16-Bit-Index, Dichte bis 50~\%), die Kompaktseite hält sortierte
  Schlüssel-Wert-Vektoren (Dichte über 50~\%); die Grenze bei 50~\% ist die Seitentyp-Grenze, die drei
  Umschaltregeln des Prüflings nennt \emph{Dichte-Schwellen}.
  \item[Rang] (\emph{rank}) Zentralitäts-Einordnung der Stand-der-Technik-Veröffentlichungen
  (Rang-1/2/3); erklärt in Abschnitt~\ref{sec:rqs}. Nicht zu verwechseln mit der Rangbildung der
  Messung (siehe \emph{Rangbildung}).
  \item[Rangbildung] (\emph{ranking}) Die Ordnung der vermessenen Permutationen nach einer Metrik
  über die drei Benchmarking-Granularitäten (Aufgabenstellung, Teilaufgabe~7); das erste Inkrement
  ist das Ranking je Metrik aus den Mess-CSVs mit Auslieferung der besten Binary
  (\texttt{best\_binary\_selector}).
  \item[RCU] Read-Copy-Update --- Reklamations-Schema (McKenney, P29)~\cite{mckenney2001rcu}.
  \item[Realm] (\emph{realm}) Einer der drei Bereiche der Achsen mit je eigener Registry,
  Registry-XML und Stempel-Zeile: der Organ-Realm (Wurzel: die Anatomie des Tiers), der
  System-Realm (Wurzel: der CEB) und der Mess-Realm (Wurzel: der Planer); die kanonische Ordnung
  lautet Mess-, System-, Organ-Achsen. Alias: Achsen-Heimat, Achsen-Realm.
  \item[Registry-XML] (\emph{registry XML}) Das je Realm per Übersetzungszeit-Reflexion generierte,
  maschinenlesbare Abbild des Achsen-Bestands (Achsen, Bausteine, Unter-Achsen,
  Maschinen-Signaturen); es wird nie von Hand editiert und ist das Angebot, gegen das der Planer
  die Anwender-XML auflöst.
  \item[REUSE, SPDX] REUSE-Spezifikation (Free Software Foundation Europe): Konvention, jede Datei
  eines Repositoriums per SPDX-Lizenzbezeichner und \texttt{REUSE.toml} einer Lizenz zuzuordnen;
  SPDX ist der standardisierte Bezeichner-Katalog für Lizenzen. Die Zuordnung lizenziert nichts um
  --- maßgeblich bleiben die Lizenzdateien~\cite{fsfe2024reuse}.
  \item[RRIP, DRRIP] (Dynamic) Re-Reference Interval Prediction ---
  Cache-Ersetzungs-Heuristik~\cite{jaleel2010rrip}.
  \item[Seitentyp] (\emph{page type}) Der Prüfling-Slot \texttt{page\_type} (prt-art
  \texttt{axis\_01}): die Knotenform des PRT-ART-B$^+$-Baums; im Kanon der Engine entspricht er der
  Organ-Haupt-Achse \texttt{node\_type} (Slot T04). Die Slot-Nummern des Prüflings folgen nicht der
  Slot-Reihenfolge T00--T17 der Engine.
  \item[SIMD] Single Instruction, Multiple Data.
  \item[SOSD] Search On Sorted Data --- Index-Benchmark~\cite{kipf2019sosd}.
  \item[SOTA] State of the Art --- der Stand der Technik.
  \item[SPD, JEDEC] Serial Presence Detect --- der von JEDEC (dem Normungsgremium der
  Speicherindustrie) standardisierte Kennungsspeicher eines Speichermoduls, der u.\,a.\ die
  JEDEC-Basisfrequenz trägt~\cite{jedec_jesd400_5}; in der Hardware-Erkennung die mittlere
  Provenienz-Stufe (SPD-JEDEC-Basiswert) zwischen konfiguriert-gemessen und
  deklariert-nicht-gemessen (siehe \emph{Provenienz-Stufe}).
  \item[Stempel] (\emph{stamp}) Die einkompilierte, ABI-stabile Identitäts-Beschriftung eines
  Trägers: je Realm eine Zeile aus Einträgen der Form \texttt{achse=algo@X.Y.Z} mit optionalen
  Flags und rekursiven Klammer-Gruppen (\texttt{c\{p.e\}} ist \texttt{c} mit den Untergliedern
  \texttt{p} und \texttt{e}, nie ein flaches Kürzel), dazu Versionszeile und Fingerprint. Vier
  Träger-Stufen und sieben Stempel-Interfaces bilden eine dreiwertige Zulassungs-Matrix von 28
  Zellen (Pflicht, Erlaubt, Verboten); der CEB trägt keine Organ-Zeile, nur der Planer eine
  Versionszeile, nur der Hybrid angeschlossene Tiere. Der Stempel ist zugleich Identität,
  Lager-Schlüssel und Wiedererkennungs-Kriterium; der Versions-Ausdruck des Treibers dient dagegen
  nur der Einordnung (Abbildung~\ref{fig:traeger-flaechen}). Alias: Identitäts-Stempel,
  Stempelsystem.
  \item[Subsystem] (\emph{subsystem}) In dieser Arbeit die vier Träger-Subsysteme der Kette
  (Mess-Treiber, CacheEngineBuilder, CacheEngine, Prüfling); im Code dagegen die zwölf
  \texttt{ISubEngine}-Stufen der Empfehlungs-Pipeline hinter der Fassade \texttt{ICacheEngine}.
  Die beiden Bedeutungen werden im Text nicht vermischt.
  \item[SuRF] Succinct Range Filter~\cite{zhang2018surf}.
  \item[SVE, NEON, RVV] ARM- (SVE, NEON) bzw.\ RISC-V- (RVV) SIMD-Erweiterungen.
  \item[Teilmenge] (\emph{subset}) In dieser Arbeit in zwei festen Bedeutungen: die zulässige
  Teilmenge der Mess-Ebenen (jede Teilmenge von wallclock/macro/micro in Kanon-Reihenfolge) und die
  Konformitäts-Teilmenge der Interface-Operationen, die das Konformitäts-Gate vor jeder Messung
  prüft (unter \texttt{--debug} zusätzlich der Genus-Testsatz; Anhang~\ref{app:interfaces}).
  \item[Tier-Binary] (\emph{tier binary}) Die dritte Träger-Stufe (Ebene~5a der Konzept-Hierarchie: die
  Nachphase der Kompilation, das Individuum): das vom CEB je Permutation
  gebaute, dynamisch ladbare Modul eines Tiers, das genau eine Organ-Belegung
  (\texttt{binary\_id}) samt System- und Mess-Wahl fest einkompiliert enthält, über buchstabengleiche
  C-Symbole (die vier Ur-Symbole, die Konzept- und Gattungs-Abfrage und die Versionszeilen) geladen
  wird und als Bestand~1 im Lager liegt. Alias: Tier, Permutations-Modul.
  \item[TLB] Translation Lookaside Buffer.
  \item[Träger-Stufe] (\emph{carrier stage}) Eine der vier Stufen, die die Anatomie tragen, in
  Bau-Reihenfolge: Planer, CacheEngineBuilder (CEB), Tier-Binary, Hybrid-Tier-Binary (jede Stufe erzeugt
  die nächste, ohne Tier kein Hybrid); in der Verarbeitungskette sitzt die Hybrid-Tier-Binary als optionaler
  Mux-Einschub zwischen CEB-Prüf-Dock und Tier-Binary. Keine
  Stufe durchläuft Schichten --- jede enthält ihre Schicht fest einkompiliert; nur die
  Konfiguration wandert von der Anwender-XML bis zur fertigen Tier-Binary. Jede Stufe führt einen
  Stempel und bietet drei Interface-Flächen (Abbildung~\ref{fig:traeger-flaechen}).
  \item[two\_phase\_valid] Gültigkeits-Flag einer Messzeile: Nur eine Messung, deren Workload über
  das zweiphasige Lastprofil (Aufbau, dann Messphase) gefahren wurde, gilt als valide; Zeilen mit
  \texttt{two\_phase\_valid} gleich \texttt{false} gehen nicht in Median, Matrix oder Diagramm ein.
%% INTERN-272: KA-150
  \item[Unter-Achse] (\emph{sub-axis}) Eine Achse unterhalb einer Haupt-Achse, im Regelfall
  dynamisch und zur Laufzeit vom Planer permutiert: Ihre Werte reisen als Einstellungen über das
  Prüf-Dock in die Tier-Binary, stehen in keiner Bau-Legende und nie in der \texttt{binary\_id}.
  Beispiele: die Laufzeit-Parameter Prefetch-Distanz, Stapelgröße, Inline-Schwelle, Pool-Budget und
  Thread-Zahl sowie die sechzehn Mess-Kategorien, der Workload und der Arbeitsmodus (Mess-Realm);
  Optimierungsstufe, Atomik-Fähigkeit, SIMD-Familie, NUMA-Knoten, Seitengröße und Kern-Klasse
  (System-Realm). Einzelne Unter-Achsen sind übersetzungsstatische Tupel ihrer Haupt-Achse (etwa
  Scheduling und Compiler) und wirken über die Bau-Version. Registry-Token \texttt{sub\_axis}.
  Alias: Sub-Achse. Unter-Achse und Permutations-Achse schließen sich nicht aus.
  \item[Verbund-Strategie] (\emph{merge strategy}) Die per Achse deklarierte Art, wie ein
  Prüfling-Slot mit der Engine-Achse verbunden wird: \texttt{Verbund1\_CeOnly} (nur die Engine,
  keine Prüflings-Direktive), \texttt{Verbund2\_Replace} (die Prüfling-Achse ersetzt die
  Engine-Achse, mit Rückfall; der Standard), \texttt{Verbund2\_Hybrid} (Engine und Prüfling hybrid
  je Prüfling) und \texttt{Verbund3\_Union} (die nicht-redundante Vereinigung); die XML-Token
  lauten \texttt{replace}, \texttt{merge} und \texttt{union}. Alias: Stufe~1 bis~3 (Altname), Full
  Join (für \texttt{union}).
  \item[Versions-Lock] (\emph{version lock}) Die eingecheckte Sperrdatei, die je Quelldatei der
  Overlay-Quellmenge Kategorie, Achsen-Version und SHA-256 verzeichnet; jede Änderung an einer
  verzeichneten Datei verlangt den bewussten Neu-Erzeugungs-Commit, eine Abweichung gegen die
  Entdeckung ist ein Fehlschlag, und ein CI-Gate prüft sie.
  \item[Virus] (\emph{virus}) Ein achsenloses Graph- oder Rein-Mathematik-Modul an derselben
  Wurzel wie die Lebewesen: ohne Achsen-System, ohne Komposition und außerhalb des
  Permutationsbaums, aber mit eigenem flachem Mess-POD vermessbar (Latenz, Durchsatz, modul-eigene
  Zähler). Plural: Viren.
  \item[Vollprofil, abstraktes Profil] (\emph{full profile, abstract profile}) Die zwei Typen der
  Stand-der-Technik-Profile: Ein Vollprofil beschreibt einen eigenständigen Suchentwurf
  vollständig; ein abstraktes Profil stellt nur einzelne Organe (OLC als
  Nebenläufigkeits-Organ~\cite{leis2016olc}, LOUDS als
  Speicher-Layout-Organ~\cite{jacobson1989louds}, VAMPIR als
  Mess-Technologie~\cite{berthold2023vampir}) und überlässt die übrigen dem Standard-Baukasten ---
  derselbe Typ, dem der Prüfling PRT-ART entspricht. Von den 33 Profilen sind 30 Vollprofile und
  3 abstrakte.
  \item[working\_set\_n] Die Spalte der Arbeitsmengen-Größe (Anzahl Schlüssel) einer Messreihe;
  sie ist die Abszisse der Sweep-Kurven des Kurvenladers, und die Zell-Token \texttt{n/a} und
  \texttt{failed} sind keine Zahlen --- ohne Messwert entsteht kein Stützpunkt.
%% INTERN-272: KA-151
  \item[XML-Lastprofil-Ergebnis] (\emph{XML load-profile result}) Das in der Aufgabenstellung
  (Teilaufgabe~8, Automatisierung der Systemerkennung und Binary-Auswahl) geforderte
  Persistenz-Format, in dem die für einen Gesamtalgorithmus extrahierte, je Architekturfokus
  wiederverwendbare Heuristik-Richtlinie unter Referenz der angewendeten Lastprofile abgelegt wird;
  dient als Eingabe der autonomen Heuristik-Profil-Auswertung. Import-Parser und Export-Writer
  samt Extraktor sind implementiert (round-trip-fähiges Schema \texttt{comdare\_load\_profile});
  Ausblick ist die darauf aufsetzende autonome Heuristik-Auswertung (ML-Klassifikation).
  \item[YCSB] Yahoo!\ Cloud Serving Benchmark~\cite{cooper2010ycsb}.
  \item[Zen~5] AMD-Mikroarchitektur (CPUID-Familie 26); die erste Messplattform dieser Arbeit
  (prod1, AMD Ryzen~9~9950X3D, Granite Ridge, zwei Kern-Komplexe) gehört ihr an, und ihr % chktex 29
  Roh-Ereignis-Katalog ist durch Messung gegengeprüft. In der System-Registry ist sie die
  Rekombination \texttt{prod1\_zen5}. Der Begriff bezeichnet in dieser Arbeit ausschließlich die
  Mikroarchitektur.
%% INTERN-272: KA-153
  \item[Zwei-Gate-Protokoll] (\emph{two-gate protocol}) Das Abnahme-Protokoll je
  Implementierungs-Paket: Gate~1 ist das Doppellauf-Gate des Engine-Repositories --- die volle
  Test-Suite läuft in vier Zellen (gcc und clang, jeweils als Release- und Debug-Bau) ohne
  Änderungs-Skip, und akzeptiert wird nur ein Lauf ohne fehlgeschlagene Tests, ohne Ausnahmeliste;
  Gate~2 ist das Integrations-Gate des Projekt-Repositories (Super-Repositorium), dessen Pipeline
  fehlerfrei ohne Ausnahmemenge baut (vgl.\ Abschnitt~\ref{sec:impl-contracts}).
\end{description}
